\section{Metodología}

\subsection{Dataset}
El Spotify Tracks Dataset contiene 114,000 canciones distribuidas en 125
géneros musicales \cite{spotifyKaggle}. Cada registro incluye metadatos
(artista, álbum, género) y 10 características de audio cuantitativas
normalizadas por Spotify mediante análisis de señal.

\subsection{Preprocesamiento}
Las características seleccionadas para el clustering fueron: danceability,
energy, loudness, speechiness, acousticness, instrumentalness, liveness,
valence, tempo y popularity. Se aplicó \texttt{StandardScaler} para
normalizar cada feature a media cero y desviación estándar unitaria,
requisito fundamental para K-Means dado su dependencia de distancias
euclidianas \cite{jain2010data}.

\subsection{Algoritmo K-Means}
K-Means minimiza la inercia intra-cluster:
\begin{equation}
J = \sum_{k=1}^{K} \sum_{\mathbf{x}_i \in C_k} \|\mathbf{x}_i - \boldsymbol{\mu}_k\|^2
\end{equation}
donde $\boldsymbol{\mu}_k$ es el centroide del cluster $C_k$.

\subsection{Selección de K}
Se evaluaron valores de $K \in \{2, \ldots, 10\}$ utilizando:
\begin{itemize}
  \item \textbf{Método del Codo}: punto de inflexión en la curva de inercia.
  \item \textbf{Silhouette Score}: mide cohesión interna y separación entre clusters.
        Rango $[-1, 1]$; valores cercanos a 1 indican clusters bien definidos.
\end{itemize}

\subsection{Reducción Dimensional --- PCA}
PCA transforma el espacio de 10 dimensiones a 2 componentes principales
para visualización, preservando la máxima varianza posible:
\begin{equation}
\mathbf{Z} = \mathbf{X} \mathbf{W}_{2}
\end{equation}
donde $\mathbf{W}_2$ contiene los dos eigenvectores de mayor eigenvalor
de la matriz de covarianza $\mathbf{\Sigma}$.
